% ============================================================
% BLOQUE 1 — Introducción (columna izquierda)
% ============================================================
\begin{block}{Introducción}
	El guiado de proyectiles es un problema central en ingeniería y defensa. 
	En este trabajo se aborda desde la \textbf{teoría de control óptimo}, 
	lo que permite derivar las leyes de guiado clásicas y extenderlas a escenarios
	más ricos: objetivos móviles, ángulos iniciales arbitrarios y coordinación 
	de múltiples proyectiles.

	\vspace{0.2cm}
	\textbf{Herramientas}: Principio de Pontryagin, ecuaciones de Riccati,
	discretización directa con \texttt{OptimalControl.jl} e Ipopt y simulación con \texttt{DifferentialEquations.jl}.

	% \insertimage{intro-foto}{width=0.7\colwidth}{Motivación: persecución target-proyectil.}
\end{block}

% ============================================================
% BLOQUE 2 — Modelo lineal (columna izquierda)
% ============================================================
\begin{block}{Modelo Lineal (PLC)}
	\begin{exampleblock}{}
	Se asume una trayectoria \textit{near-miss}: ángulos pequeños, velocidad de 
	acercamiento $V_c=-\dot R>0$ y aceleración solo en el eje perpendicular a la 
	\textit{Line of Sight} (LOS). Definiendo $z=(r_y,v_y)$, la dinámica se reduce a
	
	$$
	\dot z=\begin{pmatrix}0&1\\0&0\end{pmatrix}z
	+\begin{pmatrix}0\\-1\end{pmatrix}u,
	\qquad z(0)=\begin{pmatrix}r_{y,0}\\v_{y,0}\end{pmatrix}.
	$$
	
	El tiempo de encuentro es $t_f=-r_{x,0}/v_{x,0}$, y con 
	$t_{\mathrm{go}}=t_f-t$ el problema es
	
	$$
	\min_{u}\; J(u)=q_f\,r_y(t_f)^2+\int_0^{t_f}u^2(t)\,dt.
	$$
	
	Vía Riccati y tomando $q_f\to\infty$ se recupera la \textbf{ley PLC clásica}:
	
	$$
	\boxed{\;\vphantom{\frac{\frac{A}{B}}{\frac{C}{D}}}
	u_\infty^*(t)=\frac{3}{t_{\mathrm{go}}^2}
	\big[r_y(t)+v_y(t)\,t_{\mathrm{go}}\big]\;}
	$$
	\end{exampleblock}
	{\def\marginimagetop{-2}
		\insertimage{lineal/colision.png}{width=\colwidth}{Trayectorias del objetivo y del proyectil bajo control PLC. La estrella indica la intercepción.}
	}

	%!TODO:
	%	- Agrandar tamaño estrella
	%	- Titulo grafico
	%	- Aumentar grosor lineas 
	%	- Revisar label
	
\end{block}

\begin{block}{Modelo No Lineal}
	Se considera un objetivo fijo ($\vec v_T=\vec a_T=0$) y se utiliza la geometría exacta. En coordenadas
	$(R,L)$ con $L=\theta_M-\lambda$ la dinámica relativa es:
	
	\vspace{-0.9cm}
	$$
	\begin{aligned}
	\dot R&= -V_M\cos L, &\quad& R(0)=R_0,\quad R(t_f)=\varepsilon,\\
	\dot L&= \frac{u}{V_M}\cos L+\frac{V_M}{R}\sin L, &\quad& L(0)=L_0.
	\end{aligned}
	$$

	\vspace{1.5cm}
	
	\begin{exampleblock}{}
	Con costo $\;J(u,t_f)=\int_0^{t_f}\!\big(\tfrac{\alpha}{2}u^2+\beta\big)dt\;$
	y $t_f$ libre, \textbf{Pontryagin} da el control óptimo
	$$
	\boxed{\;\vphantom{\frac{\frac{A}{B}}{\frac{C}{D}}}
	u(t)=\operatorname{proj}_{[-u_{\max},u_{\max}]}
	\!\left(-\frac{p_\theta(t)\,\cos L}{\alpha V_M}\right)\;}
	$$
	con coestado $p_\theta$ y condición de transversalidad en $t_f$:
	$$
	\Theta(t_f)=0,\qquad p_x(t_f)=2\nu r_x(t_f),\quad p_y(t_f)=2\nu r_y(t_f),\quad p_\theta(t_f)=0.
	$$
	Si inicialmente apunta al objetivo ($L_0=0$, $R_0>\varepsilon$), la solución es $u^*\equiv0$ y
	$t_f^*=(R_0-\varepsilon)/V_M$.
	\end{exampleblock}
	{\def\marginimagetop{-2}
		\insertimage{no_lineal/comparacion.png}{width=0.85\colwidth}{Soluciones numéricas comparadas para distintos ángulos iniciales (panel derecho reflejado).}
	}
\end{block}

% ============================================================
% CREA UNA NUEVA COLUMNA
% ============================================================
\separatorcolumn

\begin{block}{Modelo de Dos Proyectiles}
	Dos proyectiles con rapidez $V_M$ buscan interceptar un objetivo en un
	\textbf{tiempo final común} $t_f$, con \textbf{ángulo relativo de llegada} $\theta_d$.
	Se aplica una aproximación lineal para cada proyectil $i\in\{1,2\}$ y se introduce
	el acoplamiento en el costo terminal, en donde se penalizan las desviaciones del ángulo de llegada:

	\vspace{-0.9cm}
	$$
	\begin{aligned}
	\min_{u_1,u_2}\;J&=\frac12\int_0^{t_f}(u_1^2+u_2^2)\,dt
	+\frac b2\sum_{i=1}^2 y_i(t_f)^2\\
	&\quad+\frac c2\big[\dot y_1(t_f)-\dot y_2(t_f)-V_M\theta_d\big]^2,
	\qquad b,c>0.
	\end{aligned}
	$$

	Al estar asumiendo el \textbf{modelo lineal} para la obtención del control óptimo, la aproximación angular pierde precisión al aumentar $\theta_d$.

	{\def\marginimagetop{-1}
		\insertimage{lineal_dos_proyectiles/comparacion_angulos_analitica.png}{width=0.9\colwidth}{Trayectorias de la solución analítica para distintos ángulos relativos deseados de llegada.}
	}
	%!TODO:
	%	- Colores
	%	- Leyenda interna
	%	- Leyenda completa
	%	- Posición theta's

\end{block}

\begin{block}{Modelo Multiproyectil}
	Se consideran un objetivo fijo $T$ y $N$ proyectiles $M_i$, $i\in I=\{1,\dots,N\}$.
	Cada proyectil usa la \textbf{aceleración normal a la velocidad} como control $u_i$, y
	sus estados $(r_{x_i},r_{y_i},\theta_{M_i})$ obedecen a la dinámica

	\vspace{-0.9cm}
	$$
	\dot r_{x_i}=-V_{M_i}\cos\theta_{M_i},\quad
	\dot r_{y_i}=-V_{M_i}\sin\theta_{M_i},\quad
	\dot\theta_{M_i}=\frac{u_i}{V_{M_i}}.
	$$
	
	Los proyectiles deben alcanzar al objetivo en un \textbf{tiempo final común} $t_f$, 
	con tolerancia $R_i(t_f)=\lVert\vec r_i(t_f)\rVert\le\varepsilon$. El acoplamiento
	aparece en las $N$ condiciones terminales simultáneas:
	
	\vspace{-0.9cm}
	$$
	h_i(x(t_f))\le0,\qquad h_i(x):=r_{x_i}^2+r_{y_i}^2-\varepsilon^2,
	\quad\forall i\in I.
	$$

	\vspace{0.5cm}
	\begin{exampleblock}{}

	Para penalizar llegadas anticipadas se añaden \textbf{holguras}
	$0\le s_i\le\varepsilon^2$ con la restricción de trayectoria
	$r_{x_i}^2+r_{y_i}^2+s_i\ge\varepsilon^2$, y el costo instantáneo pasa a ser
	
	\vspace{-0.9cm}
	$$
	\boxed{\;\vphantom{\frac{\frac{A}{B}}{\frac{C}{D}}}\;
	\ell(x,w)=\beta+\sum_{i\in I}\!\Big(\tfrac{\alpha}{2}u_i^2+\gamma s_i\Big)\;}
	$$
	El problema de Bolza a tiempo final libre queda

	$$
	\min_{\substack{t_f>0\\w=(u,s)}}\;
	J(w,t_f)=\int_0^{t_f}\!\Big(\beta+\frac{\alpha}{2}\sum_{i\in I}u_i^2
	+\gamma\sum_{i\in I}s_i\Big)dt.
	$$
	
	Al imponer la penalización de llegadas anticipadas, los proyectiles cercanos al objetivo
	se desvían de la trayectoria para que puedan alcanzarlo a la vez que los proyectiles lejanos.
	\end{exampleblock}

	{\def\marginimagetop{-2}
		\insertimage{multiproyectil/trayectorias_N8.png}{width=\colwidth}{Trayectorias de 8 proyectiles con tiempo final común y penalización de llegadas anticipadas.}
	}

\end{block}
